\section{Discussão: do impresso às redes digitais}
\subsection{O que se pode afirmar sobre a origem de \forma{kkkk}}
As fontes inspecionadas sustentam uma história de representações múltiplas do riso. Não sustentam uma linha necessária \forma{quiá} $>$ \forma{cá} $>$ \forma{quá} $>$ \forma{kakaka} $>$ \forma{kkkk}. Para testar essa linha, seriam úteis documentos intermediários, declarações contemporâneas de adaptação, padrões de circulação, reedições, tradições de quadrinhos e corpora digitais datados. Uma série de setas não substitui esses elos.

Também não se demonstra independência absoluta. Uma geração não precisa ter lido conscientemente \textit{Til} para participar de uma convenção transmitida por quadrinhos, escola, família ou mídia. Transmissão indireta é uma possibilidade; reinvenção também. O exemplo hipotético do material inicial, em que leitores adotariam uma forma dos livros e a passariam adiante, é uma hipótese legítima. Seu requisito não deve ser estreitado a fala cotidiana nem a manutenção de um radical identificável.

A formulação sustentada é: há antecedentes impressos funcional e parcialmente formalmente semelhantes ao riso digital; a coleta não estabelece uma cadeia específica de descendência, tampouco sua inexistência. A forma digital está documentada no conjunto acadêmico ligado a 2007.2 e pode ser anterior. O intervalo ``1998 a 2002'' não foi validado como origem. Atribuí-lo a consenso de historiadores ou linguistas, sem estudos identificados, seria um argumento de autoridade sem referência.

\subsection{IRC, mIRC, BBS, salas públicas e mensageiros}
O material enumera mIRC, ICQ, bate-papos de UOL, BOL e Terra, MSN, Orkut e jogos. Esses ambientes são relevantes para investigar a difusão de convenções digitais, mas precisam ser diferenciados. IRC é um protocolo ou ambiente de conversação; mIRC é um cliente. Mensageiros privados, fóruns, serviços de recados e salas públicas têm formas distintas de arquivamento, público e temporalidade. ``Forjado sequencialmente em três plataformas'' exige evidência de transferência entre elas.

O levantamento não recuperou logs autenticados de 1995--1997 que sustentem a alegação de ausência, nem uma coleção de 1999--2000 que demonstre \forma{kakaka} como estágio obrigatório. A existência histórica dessas plataformas não comprova a presença de uma grafia específica. Um testemunho retrospectivo pode orientar buscas, mas está sujeito a memória reconstrutiva. Sua data é a do depoimento, salvo quando acompanhado de arquivo contemporâneo verificável.

As alegações de massificação pelo MSN e de uso em jogos como \textit{Ultima Online}, \textit{Tibia} e \textit{Counter-Strike} permanecem hipóteses de difusão. A afirmação de banimentos frequentes por confusão com uma sigla supremacista precisaria de regras de moderação, mensagens administrativas ou registros datados. Sem esses documentos, não se transforma uma anedota plausível em prova de origem.

\subsection{A hipótese ergonômica e seus limites causais}
Repetir uma tecla pode ser mecanicamente mais simples que alternar duas. O nome \forma{cá} também oferece uma interpretação possível da cadeia de letras. Isso torna plausível uma hipótese de economia gráfica, mas não prova a motivação dos primeiros usuários. A escrita digital contém inúmeras formas mais longas e custosas, como \forma{huahuahua} e sequências heterogêneas, empregadas por seus efeitos expressivos. A redução de esforço é um fator possível entre outros.

A hipótese poderia ser testada experimentalmente com tempos de digitação, familiaridade, dispositivo e tarefa controlados, ou documentalmente com entrevistas contemporâneas e séries de uso. Uma experiência atual demonstraria diferenças atuais de custo, não a causa histórica original. Também se deveria distinguir repetição por toque, pressionamento contínuo e recursos de autocompletar: diferentes sistemas e teclados comportam-se de maneiras distintas.

A presença ou ausência de \forma{k} no inventário alfabético normativo não determina sua disponibilidade em máquinas ou na escrita. Letras podem circular em nomes, símbolos e empréstimos antes de mudanças normativas. A reforma ortográfica não é uma evidência de datação de \forma{kkkk}. O argumento de que ``a letra só entrou no alfabeto depois'' confunde norma ortográfica, práticas sociais e existência material do caractere.

\subsection{Tecnologia, imprensa e coexistência de convenções}
Tipografia, composição editorial, espaço de coluna e práticas de tradução podem influenciar o texto. Entretanto, não se deve atribuir a cada década uma personalidade coletiva: riso ``aristocrático'' em 1900, ``malandro'' em 1960 e ``funcional'' em 1980 são interpretações que precisariam de indicadores. Uma passagem de folhetim, uma rubrica parlamentar e um balão de quadrinhos pertencem a gêneros distintos, mesmo quando publicados no mesmo ano.

A explicação de que a censura dos anos 1940 ou 1970 substituiu onomatopeias por rubricas exigiria comparar publicações, editorias e períodos sob procedimentos definidos. Correlação temporal com um regime político não estabelece causalidade gráfica. Do mesmo modo, a popularidade da Disney não demonstra que o \forma{quá} brasileiro se originou ali; a capa de 1935 fornece um testemunho anterior ao marco proposto. A força de uma indústria cultural é compatível com difusão de formas já existentes.

As hipóteses sobre origem de \forma{rs} em telegramas, gabinetes e redações merecem investigação, não endosso automático. Seriam necessários manuais, cartas, rascunhos ou laudas datadas com interpretação inequívoca. A facilidade de explicar uma abreviação a partir de \forma{risos} não fornece sua data inicial. Até que se recupere esse material, a genealogia específica ``máquina de escrever $>$ chat'' permanece indeterminada.

\section{Variação pragmática contemporânea e o problema das gerações}
\subsection{Maiúsculas, extensão e sequências heterogêneas}
\forma{kkkk}, \forma{KKKKKK}, \forma{hehe}, \forma{hahaha}, \forma{rsrs} e \emph{keysmashes} podem diferenciar intensidade, afiliação, ironia, constrangimento ou estilo. A interpretação depende de posição na conversa, relação entre interlocutores, gênero discursivo, plataforma e repertório local. Maiúsculas e alongamento podem intensificar a expressão, mas não há fundamento para uma regra universal segundo a qual minúsculas indiquem falsidade e maiúsculas sinceridade.

Uma cadeia heterogênea como \forma{asdfghjkl} pode figurar descontrole expressivo e funcionar como prática aprendida. Não precisa ter surgido com telas de toque: sequências baseadas no alinhamento das teclas são compatíveis com teclado físico. Um formato pode ser empregado em diferentes dispositivos, e uma explicação causal precisa datar sua ocorrência. Não se infere o aparelho utilizado apenas pela forma textual.

\subsection{Emoticons, emojis e respostas verbais}
Emoticons como \texttt{:-)} e emojis são recursos distintos. Os primeiros combinam caracteres de escrita; os segundos são sinais pictográficos codificados, com desenho que pode variar entre plataformas. Emoji de caveira, rosto chorando e rosto com lágrimas de alegria podem adquirir leituras humorísticas, mas o valor convencional não se esgota numa tradução fixa. O uso de caveira para humor pode explorar a expressão ``morrer de rir''; isso não prova que toda ocorrência de caveira denote gargalhada.

O repertório também inclui respostas como \forma{rindo}, \forma{rindo de nervoso}, comentários lacônicos, imagens e figurinos visuais. Uma resposta com ponto final pode comunicar secura ou ironia em determinado grupo, sem que o ponto carregue sempre esse significado. A revisão de Badrakh e colaboradores \cite{badrakh2024} destaca a relevância e a ambiguidade dos sinais digitais em brincadeiras online. A interpretação precisa da sequência conversacional; uma lista de equivalências descontextualizadas perde esse componente.

O significado de um marcador pode mudar sem alteração em sua forma. Esse é um caso pertinente de mudança semântica ou pragmática. Portanto, o fato de não haver derivação de um radical não exclui mudança linguística. A expansão de leituras humorísticas de um signo visual permite precisamente investigar convencionalização, inferência e circulação histórica.

\subsection{Geração Z e Alpha: hipóteses, não gramáticas uniformes}
As respostas iniciais atribuem à geração Z e à Alpha preferências uniformes, com limites de nascimento fixos e oposição rígida aos millennials. O presente corpus não contém amostra brasileira por idade que permita estimar tais preferências, e não inclui dados próprios sobre crianças da geração Alpha. Identificar uma tendência num grupo online não autoriza generalização a todos os integrantes de uma coorte.

Kempe e Raviv \cite{kempe2025} estudam a convencionalização de emojis faciais e não encontram as diferenças geracionais previstas naquele desenho. O resultado não elimina toda variação etária, cultural ou contextual, mas impede tratar manchetes sobre ``emoji de velho'' como regra científica universal. A revisão mais ampla de Costanza, Rudolph e Zacher \cite{costanza2023} também alerta para a fragilidade de atribuições geracionais que confundem idade, período e coorte.

Uma investigação brasileira precisaria recrutar participantes com idades conhecidas, caracterizar região e comunidades de prática, medir usos e interpretações em tarefas comparáveis e analisar diferenças de plataforma. Para menores de idade, seriam necessários procedimentos éticos apropriados ao desenho. O corpus documental deste artigo não equivale a uma pesquisa com participantes, e não se declara aprovação ética inexistente.

\section{Bibliografia, autoridade e integridade da evidência}
\subsection{Referências da UnB e atualização teórica}
A escolha de bibliografia não se faz apenas pela data mais recente. Obras de referência podem continuar relevantes; artigos recentes são indispensáveis para debates e dados novos. O Edital 04/2026 da UnB fornece uma base institucional verificável para Martelotta e Seara, Nunes e Lazzarotto-Volcão \cite{unb2026}. A bibliografia inclui ainda estudos recentes sobre discurso e multimodalidade. Não se presume que a lista seletiva esgote o ensino ou a pesquisa do departamento.

Campbell \cite{campbell2021} oferece enquadramento histórico-comparativo contemporâneo; Ilari \cite{ilari2018}, uma referência para a formação românica. Monte \cite{monte2019} é especialmente importante porque associa discussão de tratamento a documentação datada. Iida e Akita \cite{iida2024} atualizam a relação entre experiência e iconicidade. Os estudos sobre sinais digitais e emojis ampliam a interpretação pragmática sem fornecer, indevidamente, a origem do \forma{kkkk} brasileiro.

\subsection{Revisão das referências sugeridas no material inicial}
Saussure, Coseriu, Ullmann, Bréal, Jakobson, Tarallo, Faria, Christiansen e Chater são nomes associados a tradições relevantes. Contudo, sua autoridade não autentica uma tese específica sem passagem pertinente. Não se atribui a Guimarães a demonstração de que onomatopeias não possuem raízes históricas, nem a Jakobson uma regra de independência de todas as imitações sonoras. Títulos, edições e usos argumentativos precisam ser conferidos antes da inclusão como sustentação.

O material atribui a Ilari um livro denominado \textit{Linguística histórica}, com dados bibliográficos não corroborados na coleta; a obra identificada é \textit{Linguística românica}. Esse tipo de correção não é cosmético: permite que outro leitor encontre a fonte. Também é inadequado transformar um livro sobre fonética do latim em autoridade suficiente para a história de qualquer palavra portuguesa sem localizar a discussão correspondente.

Uma referência acadêmica recente pode repetir uma hipótese de divulgação sem fornecer nova prova primária. Identificou-se, por exemplo, um artigo de 2025 que remete a uma página de divulgação ao narrar a relação entre \forma{quiá} e \forma{kkkk}. Esse encadeamento mostra circulação da narrativa, não confirmação independente da genealogia. Contar citações dessa versão como testemunhos distintos produziria uma falsa multiplicação de evidências.

\subsection{Fidelidade ao material de partida e transparência de autoria}
Ser fiel às questões enviadas exige preservar o seu problema, suas alternativas e seus exemplos, mas não manter como verdade o que a verificação não sustenta. Por isso, o artigo conserva a discussão sobre raiz, mudança, onomatopeia, suporte, periódicos, plataformas e gerações; reclassifica as supostas provas; corrige conceitos e registra lacunas. As citações sem autenticação aparecem como alvos de auditoria, e não como falas efetivas dos autores históricos.

Murillo Edson de Carvalho Souza é o único autor identificado. A elaboração deste manuscrito e dos scripts utilizou assistência de inteligência artificial para organização, redação, programação e busca documental. Ferramentas automatizadas não são coautoras nem fontes históricas. A publicação do material não equivale a revisão humana independente ou avaliação por pares; a responsabilidade pela versão científica submetida e por suas conclusões permanece com o autor. O artigo não herda do template as declarações, afiliações, dados ou coautoria do texto que nele vinha exemplificado.
